%% FT_QuestionPool_Additional_MC.tex
%% ─────────────────────────────────────────────────────────────────────────────
%% ADDITIONAL CONTROL QUESTIONS (Single Choice / type=mc)
%% Foreign Trade and International Organisations, Jan Suchanek, Hochschule
%% Bremen, Summer Semester 2026
%%
%% This file is a SUPPLEMENT to FT_QuestionPool.tex. It contains only new
%% \begin{question}[type=mc]{...}...\end{question} blocks, grouped under the
%% same chapter headings as the master file, so each block can be copy-pasted
%% straight into the matching \section{...} of the main question pool.
%%
%% New material drawn from the script that was not yet covered by the existing
%% pool: cross-price elasticity, the fundamental equation of macroeconomics,
%% and trade-openness indicators (Ch.1); TEU/FCL-LCL, VAT vs. excise, and
%% Rana Plaza/CSDDD (Ch.2); Dutch Disease, the rentier state, and the elephant
%% curve (Ch.3); the gravity model of trade and Thünen's rings (Ch.4); the
%% median voter theorem and the branches/levels of the state (Ch.5); embargoes
%% & sanctions and tariff retaliation (Ch.6); the concentration ratio and
%% Porter's Six Forces (Ch.7); UNCLOS/exclusive economic zones and the
%% freedom of the seas (Ch.8); the Phillips curve/NAIRU and the US dollar's
%% reserve-currency role (Ch.9); purchasing power parity, the Marshall-Lerner
%% condition, and the Mundell-Fleming model (Ch.10).
%%
%% To compile this file standalone for proofreading, use the same preamble as
%% FT_QuestionPool.tex (hsbquestions.sty, hsbcommon.sty must be present).
%% ─────────────────────────────────────────────────────────────────────────────

\documentclass[paper=a5]{scrartcl}
\usepackage{xcolor}
\usepackage{hsbcommon}
\usepackage[hidelinks]{hyperref}
\usepackage[utf8]{inputenc}
\usepackage{setspace}
\usepackage{pgfplots}
\pgfplotsset{compat=1.18}
\usepackage{hsbquestions}
\showanswerstrue                 %% ANSWER KEY view for proofreading — set to
                                  %% \showanswersfalse before distributing.

\title{Foreign Trade and International Organisations}
\subtitle{Additional Control Questions (Single Choice) — Supplement}
\author{Claude (Anthropic) using the Script of Jan \noun{Suchanek}}

\begin{document}
\onehalfspacing
\maketitle
\clearpage

%% ═════════════════════════════════════════════════════════════════════════════
\section{Chapter 1 — The World Economy and National Accounts}
%% ═════════════════════════════════════════════════════════════════════════════

\begin{question}[type=mc]{Cross-Price Elasticity of Demand}
The cross-price elasticity of demand between two goods $x$ and $y$ is
defined as $XED_{x,y} = \frac{dx/x}{dp_y/p_y}$. A \emph{negative} value of
$XED_{x,y}$ indicates that:
\begin{choices}
  \choice        $x$ and $y$ are complements.
  \correctchoice $x$ and $y$ are substitutes.
  \choice        $x$ is a Giffen good.
  \choice        $y$ is a public good.
\end{choices}
\begin{solution}
  If the price of $y$ rises and the quantity demanded of $x$ \emph{falls}
  ($XED_{x,y} < 0$), the two goods must be \textbf{complements} --- they tend
  to be consumed together, so less $y$ means less demand for $x$ too. If a
  price rise in $y$ instead \emph{increases} demand for $x$ ($XED_{x,y}>0$),
  the goods are \textbf{substitutes}: buyers switch away from the
  now-more-expensive $y$ toward $x$. Examples in the script: copper and
  aluminium are substitutes; copper and copper wire are complements.
\end{solution}
\end{question}

\begin{question}[type=mc]{The Fundamental Equation of Macroeconomics}
The expenditure approach to GDP is expressed as $Y = C + I + Ex - Im + G$.
Which of the following correctly interprets this identity?
\begin{choices}
  \choice        Imports must always equal exports for the equation to hold.
  \correctchoice Everything produced in a country ($Y$) is either consumed,
                 invested, exported, or represents state expenditure, net of
                 what was imported rather than produced domestically.
  \choice        Government spending $G$ is a subset of investment $I$.
  \choice        The equation only holds for closed economies.
\end{choices}
\begin{solution}
  This account-based identity says that everything produced domestically
  ($Y$) is used up somewhere: it is \textbf{consumed} ($C$), \textbf{invested}
  ($I$, which equals savings $S+\Delta K$), \textbf{exported} ($Ex$), or
  represents \textbf{government expenditure} ($G$); \textbf{imports} ($Im$)
  are subtracted because they are demand satisfied by goods \emph{not}
  produced domestically. It holds for any open economy, not just the closed
  (world) economy, precisely because $Ex$ and $Im$ are included as separate
  terms.
\end{solution}
\end{question}

\begin{question}[type=mc]{Trade Openness, Import Penetration, and Dependency}
A country's \emph{Trade Openness} is calculated as $(|EX|+|IM|)/Y$, while
its \emph{Import Penetration Rate} for a specific good is $IM/Y$. What is
the key limitation of the Trade Openness measure that the Import Penetration
and Dependency Rate try to address?
\begin{choices}
  \choice        Trade Openness cannot be calculated for small countries.
  \correctchoice Trade Openness treats all trade alike and double-counts
                 cross-border flows, without revealing what specific goods
                 are traded or how dependent the country is on particular
                 imports.
  \choice        Trade Openness always equals exactly 100\,\% for large
                 economies.
  \choice        Import Penetration measures only the volume of exports, not
                 imports.
\end{choices}
\begin{solution}
  \textbf{Trade Openness} ($TO = (|EX|+|IM|)/Y$) counts every cross-border
  transaction twice (once as an export somewhere, once as an import
  elsewhere) and says nothing about \emph{what} is traded or with whom.
  Large economies are naturally ``less open'' than small ones simply because
  a bigger share of their output is consumed domestically --- this reveals
  nothing about trade policy. The \textbf{Import Penetration Rate} and the
  good-specific \textbf{Dependency Rate} ($DR_X = IM_X/C_X$) instead show how
  reliant a country is on imports of a particular good (e.g.\ energy or a
  critical raw material).
\end{solution}
\end{question}


%% ═════════════════════════════════════════════════════════════════════════════
\section{Chapter 2 — Freight and Transport}
%% ═════════════════════════════════════════════════════════════════════════════

\begin{question}[type=mc]{TEU and Container Loads}
In shipping terminology, ``TEU'' and the distinction between ``FCL'' and
``LCL'' refer to:
\begin{choices}
  \choice        TEU is a customs code; FCL/LCL describe insurance categories.
  \correctchoice TEU (Twenty-foot Equivalent Unit) is the standard unit for
                 measuring container-ship capacity; FCL (Full Container Load)
                 means one shipper fills an entire container, while LCL (Less
                 than Container Load) means several shippers' goods are
                 grouped (groupage) into one container.
  \choice        TEU measures deadweight tonnage; FCL/LCL describe whether a
                 ship is a liner or a tramp vessel.
  \choice        TEU is only used for bulk liquid cargo such as LNG.
\end{choices}
\begin{solution}
  \textbf{TEU} (Twenty-foot Equivalent Unit) standardises container-ship
  capacity so vessels of different sizes can be compared. \textbf{FCL} means
  a single consignor's cargo fills the whole container; \textbf{LCL} means
  multiple smaller shipments are consolidated (groupage) into a shared
  container. \textbf{DWT} (deadweight tonnage) is a separate measure of how
  much weight a ship can carry.
\end{solution}
\end{question}

\begin{question}[type=mc]{VAT vs.\ Excise Duty}
When goods cross a border, they may be subject to both VAT and an excise
duty. What is the key difference between the two?
\begin{choices}
  \choice        VAT applies only to services, excise duty only to goods.
  \correctchoice VAT is a general tax levied on the value of (almost) all
                 goods, while excise duties are special taxes levied only on
                 particular goods, such as tobacco, alcohol, or fuel.
  \choice        Excise duty is refundable, while VAT is never refundable.
  \choice        VAT is levied only within the EU, excise duty only outside
                 the EU.
\end{choices}
\begin{solution}
  \textbf{VAT} (value added tax) is a broad-based tax on the value of
  virtually all goods and services. \textbf{Excise duties}
  (\emph{Verbrauchssteuer}) are narrow, targeted taxes on specific goods such
  as tobacco, alcohol, and mineral oil. VAT already paid in the exporting
  country may sometimes be deducted, whereas the buyer may still owe VAT in
  the importing country --- whether VAT is due at the border is a
  complicated, jurisdiction-specific question.
\end{solution}
\end{question}

\begin{question}[type=mc]{Rana Plaza and Corporate Due Diligence}
The 2013 collapse of the Rana Plaza textile factory in Bangladesh, in which
over 1,000 workers died, is used in the script to illustrate:
\begin{choices}
  \choice        The success of self-regulation by multinational retailers.
  \correctchoice How international trade separates consumers from producers,
                 motivating supply-chain due diligence laws such as the
                 German \emph{Lieferkettensorgfaltspflichtengesetz} (LkSG)
                 and the EU's CSDDD.
  \choice        The advantages of vertical integration in the textile
                 industry.
  \choice        A case in which the WTO imposed binding sanctions on
                 Bangladesh.
\end{choices}
\begin{solution}
  Because global trade separates the point of consumption from the point of
  production, buyers often have little visibility into working conditions
  and safety standards deep in their supply chain. Rana Plaza (like the 1984
  Bhopal disaster) galvanised support for \textbf{supply-chain due diligence}
  legislation: Germany's \textbf{LkSG} (applying to companies with over
  1,000 employees), the EU's \textbf{Corporate Sustainability Due Diligence
  Directive (CSDDD)}, and the \textbf{Corporate Sustainability Reporting
  Directive (CSRD)}.
\end{solution}
\end{question}


%% ═════════════════════════════════════════════════════════════════════════════
\section{Chapter 3 — Terms of Trade and the Effects of Trade}
%% ═════════════════════════════════════════════════════════════════════════════

\begin{question}[type=mc]{Dutch Disease}
``Dutch Disease'' describes a situation in which a natural-resource boom
(e.g.\ a new gas field) harms the rest of the economy because:
\begin{choices}
  \choice        The government is forced to nationalise the resource sector.
  \correctchoice Resource export revenues appreciate the domestic currency,
                 making other (``lagging'') export sectors less competitive
                 internationally.
  \choice        Resource extraction always causes hyperinflation within a
                 year.
  \choice        It refers to the depletion of the resource itself, not to
                 any currency effect.
\end{choices}
\begin{solution}
  A resource boom brings in large foreign-currency revenues, which
  \textbf{appreciates the domestic currency}. This makes the country's other
  exporting industries (manufacturing, agriculture) less price-competitive
  abroad --- the ``lagging sector''. The term originates from the
  Netherlands' experience after discovering North Sea gas, and it can also
  complicate inflation control and exchange-rate stabilisation.
\end{solution}
\end{question}

\begin{question}[type=mc]{The Rentier State}
A ``rentier state'' is characterised by:
\begin{choices}
  \choice        Heavy reliance on progressive income taxation.
  \correctchoice Most government revenue coming from a single sector (such as
                 mining or oil), largely independent of broader domestic
                 productivity, often with limited democratic accountability.
  \choice        A government that primarily rents out public infrastructure
                 to foreign companies.
  \choice        An economy with no natural resources at all.
\end{choices}
\begin{solution}
  In a \textbf{rentier state}, the government draws most of its revenue from
  a single resource sector rather than from taxing the productivity of its
  citizens. Because this income does not depend on broad-based economic
  performance, there is less incentive to invest in productivity growth, and
  such states --- often petrostates --- tend to exhibit weaker democratic
  participation, since citizens have less leverage as taxpayers.
\end{solution}
\end{question}

\begin{question}[type=mc]{The Elephant Curve}
The ``elephant curve'' (Lakner-Milanovic curve), plotting global income
growth by income percentile, is best known for showing that:
\begin{choices}
  \choice        Only the poorest 10\,\% of the world's population benefited
                 from globalisation.
  \correctchoice Nearly everyone's income grew, but the global upper-middle
                 class gained the most, while a segment of the (formerly)
                 rich in high-income countries saw stagnant income growth.
  \choice        Global inequality, as measured by the Gini coefficient,
                 fell to zero.
  \choice        Income growth was identical across all percentiles of the
                 global distribution.
\end{choices}
\begin{solution}
  The \textbf{elephant curve} shows real income growth across the global
  income distribution over recent decades. Almost everyone gained, but the
  \textbf{global middle} (largely emerging-market populations) gained the
  most proportionally, while a group just below the very top --- often
  working- and middle-class households in high-income countries --- saw
  their incomes stagnate. The very top percentile gained substantially again,
  producing the curve's distinctive ``elephant'' shape.
\end{solution}
\end{question}

\begin{question}[type=mc]{Non-Tradable Goods}
Why is a haircut typically far cheaper in a low-wage country than in a
high-wage country, even though a laptop costs roughly the same everywhere?
\begin{choices}
  \choice        Laptops are non-tradable, while haircuts are tradable.
  \correctchoice A haircut is a non-tradable service tied to local labour
                 costs, whereas a laptop is a tradable good whose price is
                 set on world markets and arbitraged across borders.
  \choice        Haircuts are subject to much higher import tariffs than
                 laptops.
  \choice        Labour is highly mobile across borders, which equalises
                 wages worldwide.
\end{choices}
\begin{solution}
  \textbf{Tradable goods} (laptops, commodities) can be shipped and
  arbitraged, so the Law of One Price pulls their prices together worldwide.
  \textbf{Non-tradable} services (haircuts, most personal services) cannot
  be exported or imported, so their price simply reflects local labour
  costs. Since international labour migration is small relative to the
  trade of goods, wage --- and hence non-tradable price --- differences
  across countries persist.
\end{solution}
\end{question}


%% ═════════════════════════════════════════════════════════════════════════════
\section{Chapter 4 — Trade Theory}
%% ═════════════════════════════════════════════════════════════════════════════

\begin{question}[type=mc]{The Gravity Model of Trade}
The gravity model of trade predicts the trade volume $T_{ij}$ between two
countries as $T_{ij} = A \cdot Y_i Y_j / D_{ij}$. This implies that, all else
equal:
\begin{choices}
  \choice        Trade volume is unaffected by the size of either economy.
  \correctchoice Trade volume rises with the economic size (GDP) of both
                 trading partners and falls with the distance between them.
  \choice        Only geographic distance matters; GDP has no effect.
  \choice        Trade volume is inversely proportional to both countries'
                 GDP.
  
\end{choices}
\begin{solution}
  Borrowed from Newton's law of gravity, the model predicts trade volume
  rises with the product of the two countries' \textbf{GDPs} ($Y_i \cdot
  Y_j$) --- bigger economies trade more --- and falls with \textbf{distance}
  ($D_{ij}$), which proxies for transport costs, transaction costs, risk,
  and trade barriers between them. It is one of the most empirically
  successful models in international economics.
\end{solution}
\end{question}

\begin{question}[type=mc]{Thünen's Rings}
Johann Heinrich von Thünen (1826) observed that agricultural land use forms
concentric rings around a town. What did he identify as the main driving
force behind this pattern?
\begin{choices}
  \choice        Differences in soil fertility between rings.
  \correctchoice Transportation costs, which made perishable or bulky goods
                 uneconomical to grow far from the market town.
  \choice        Government zoning regulations.
  \choice        Climate differences within a small radius of a town.
\end{choices}
\begin{solution}
  Thünen explained the rings (market gardens, then firewood, then arable
  farming, then livestock, moving outward) through \textbf{transport costs}:
  in his era, perishable or heavy goods were expensive to move by cart over
  long distances, so land uses with high transport costs per unit value were
  located closest to the town. This laid the foundation for modern models of
  economic geography, though modern transport modes (cheap sea freight, fast
  air freight) can weaken the pattern today --- fresh flowers are now often
  air-freighted from entirely different climate zones.
\end{solution}
\end{question}


%% ═════════════════════════════════════════════════════════════════════════════
\section{Chapter 5 — The State, Taxes and Tariffs}
%% ═════════════════════════════════════════════════════════════════════════════

\begin{question}[type=mc]{The Median Voter Theorem}
The Median Voter Theorem states that, in a majority-rule election along a
single left-right political spectrum:
\begin{choices}
  \choice        The most extreme candidate always wins.
  \correctchoice Any winning majority must include the voter positioned
                 exactly in the middle of the distribution of opinion --- the
                 median voter.
  \choice        Elections always end in a 50:50 tie.
  \choice        Voter turnout has no effect on the outcome.
\end{choices}
\begin{solution}
  If political opinion can be ordered along a single axis, any coalition
  that wins a majority must, by definition, include the voter standing
  exactly at the 50th percentile --- the \textbf{median voter}. This gives
  candidates a strong incentive to position their platforms near the centre
  of the distribution of opinion to capture that decisive vote, rather than
  catering only to their core supporters.
\end{solution}
\end{question}

\begin{question}[type=mc]{Branches and Levels of the State}
A typical modern state has three institutional branches. These are:
\begin{choices}
  \choice        Federal, regional (L\"ander), and communal government.
  \correctchoice The executive, the legislative, and the judicative branch.
  \choice        The central bank, the treasury, and the customs authority.
  \choice        Public law, private law, and commercial law.
\end{choices}
\begin{solution}
  The three classic branches are the \textbf{executive} (the government and
  administration), the \textbf{legislative} (parliament), and the
  \textbf{judicative} (the courts). These are conceptually distinct from the
  \emph{levels} of government (federal, regional/Länder, communal, social
  security, supranational/EU) and from the central bank, which sits outside
  this three-branch structure altogether.
\end{solution}
\end{question}


%% ═════════════════════════════════════════════════════════════════════════════
\section{Chapter 6 — Subsidies in an Open Economy}
%% ═════════════════════════════════════════════════════════════════════════════

\begin{question}[type=mc]{Embargoes and Sanctions}
Which statement about trade embargoes and sanctions is correct?
\begin{choices}
  \choice        Sanctions can only target entire countries, never specific
                 individuals or organisations.
  \correctchoice Sanctions and embargoes can target countries, specific
                 goods, or specific persons/organisations, and typically
                 redirect trade toward third, neutral countries rather than
                 eliminating it entirely.
  \choice        Unilateral sanctions have been shown to have no measurable
                 effect on the GDP of the targeted country.
  \choice        All sanctions require United Nations Security Council
                 authorisation to be legally valid.
\end{choices}
\begin{solution}
  Sanctions regimes vary in scope: they may target \textbf{whole countries},
  \textbf{specific goods} (e.g.\ arms embargoes), or \textbf{named persons
  and organisations} (asset freezes, travel bans). Rather than eliminating
  trade, sanctions typically \textbf{redirect trade flows} toward
  third/neutral countries. Unlike UN-mandated sanctions, unilateral coercive
  measures do not require Security Council authorisation, though research
  cited in the script finds they do tend to reduce the targeted country's
  GDP growth.
\end{solution}
\end{question}

\begin{question}[type=mc]{Retaliation and the Optimum Tariff}
The script contrasts the case of a large country imposing a tariff with that
of a small country. Why might a \emph{large} country consider an
``optimum tariff'' even though tariffs are generally inefficient?
\begin{choices}
  \choice        Large countries are exempt from WTO tariff-binding
                 commitments.
  \correctchoice A large country can influence the world price of the traded
                 good, so a tariff can shift part of the burden onto foreign
                 exporters --- an option unavailable to a small,
                 price-taking country.
  \choice        Large countries never face retaliation from trading
                 partners.
  \choice        An optimum tariff eliminates deadweight loss entirely.
\end{choices}
\begin{solution}
  A \textbf{small} country is a price-taker: its tariff cannot move the
  world price, so the entire burden falls domestically. A sufficiently
  \textbf{large} country, however, can affect world demand for the good
  enough that part of the tariff's cost is effectively shifted onto foreign
  exporters (a terms-of-trade gain) --- the logic behind the theoretical
  ``optimum tariff''. This does not eliminate deadweight loss, and it
  invites \textbf{retaliation}, which can erase the gain in a trade war.
\end{solution}
\end{question}


%% ═════════════════════════════════════════════════════════════════════════════
\section{Chapter 7 — Multinational Enterprises}
%% ═════════════════════════════════════════════════════════════════════════════

\begin{question}[type=mc]{Concentration Ratio vs.\ HHI}
The Concentration Ratio $CR_n = \sum_{i=1}^n a_i$ (the combined market share
of the top $n$ firms) differs from the Herfindahl-Hirschman Index mainly in
that:
\begin{choices}
  \choice        $CR_n$ can exceed 100\,\%, while HHI cannot.
  \correctchoice $CR_n$ only looks at the largest few firms and ignores the
                 structure of the rest of the market, whereas the HHI
                 (summing squared market shares of all firms) is considered
                 scientifically superior because it reflects the entire
                 distribution of market shares.
  \choice        HHI can only be computed for regulated industries.
  \choice        $CR_n$ is used only in the EU, HHI only in the US.
\end{choices}
\begin{solution}
  $CR_n$ simply adds up the market shares of the top $n$ firms (e.g.\
  $CR_4$), giving a quick but crude picture. The \textbf{HHI}, $\sum_i
  a_i^2$, weights larger firms more heavily (because shares are squared) and
  incorporates the shares of \emph{all} firms in the market (often computed
  for $N=50$ for practical purposes), which is why competition authorities
  regard it as the scientifically superior measure of concentration.
\end{solution}
\end{question}

\begin{question}[type=mc]{Porter's Six Forces}
Porter's original ``Five Forces'' framework for industry analysis is
sometimes extended to ``Six Forces''. What is typically added as the sixth
force?
\begin{choices}
  \choice        Government regulation.
  \correctchoice Complementors --- producers of complementary goods or
                 services that increase demand for the industry's product.
  \choice        Currency risk.
  \choice        Corporate social responsibility.
\end{choices}
\begin{solution}
  Porter's classic Five Forces are: rivalry among existing competitors, the
  threat of new entrants, the threat of substitutes, and the bargaining
  power of suppliers and of buyers. The commonly added sixth force is
  \textbf{complementors}: firms whose products increase demand for one's
  own (e.g.\ app developers for a smartphone platform), recognising that
  some players make the pie bigger rather than fighting over the existing
  slice.
\end{solution}
\end{question}


%% ═════════════════════════════════════════════════════════════════════════════
\section{Chapter 8 — International Law and Governance}
%% ═════════════════════════════════════════════════════════════════════════════

\begin{question}[type=mc]{Exclusive Economic Zones under UNCLOS}
Under the United Nations Convention on the Law of the Sea (UNCLOS), a
coastal state's Exclusive Economic Zone (EEZ) extends up to:
\begin{choices}
  \choice        12 nautical miles from the coastline.
  \choice        24 nautical miles from the coastline (the contiguous zone).
  \correctchoice 200 nautical miles from the coastline (or to the end of the
                 continental shelf, if further).
  \choice        There is no fixed distance; it is set bilaterally for every
                 pair of neighbouring states.
\end{choices}
\begin{solution}
  UNCLOS defines several maritime zones: \textbf{territorial waters} extend
  12 nautical miles (full sovereignty), the \textbf{contiguous zone} extends
  to 24 nautical miles (limited enforcement rights), and the \textbf{EEZ}
  extends up to \textbf{200 nautical miles} --- or to the edge of the
  continental shelf if that extends further --- within which the coastal
  state has exclusive rights to natural resources (fishing, seabed minerals,
  pipelines, cables). Disputes over EEZ boundaries are common, especially
  around small islands that generate outsized maritime claims.
\end{solution}
\end{question}

\begin{question}[type=mc]{Freedom of the Seas}
The doctrine of ``mare liberum'' (freedom of the seas), first articulated by
Hugo Grotius in the 17th century, holds that:
\begin{choices}
  \choice        All coastal states may claim unlimited sovereignty over
                 adjacent waters.
  \correctchoice The high seas are open to all nations for navigation, and no
                 state may claim sovereignty over them.
  \choice        Ships must always sail under the flag of the country where
                 their cargo originated.
  \choice        Only signatories to UNCLOS may navigate international
                 waters.
  
\end{choices}
\begin{solution}
  Grotius's \textbf{mare liberum} argued that the open sea cannot be owned
  by any single nation and must remain free for navigation and trade by all.
  This principle, combined with the later convention that a ship is governed
  by the law of the state whose flag it flies, underpins the modern law of
  the sea --- including UNCLOS, which formalised territorial waters,
  contiguous zones, and exclusive economic zones while preserving freedom of
  navigation beyond them.
\end{solution}
\end{question}


%% ═════════════════════════════════════════════════════════════════════════════
\section{Chapter 9 — Monetary Theory}
%% ═════════════════════════════════════════════════════════════════════════════

\begin{question}[type=mc]{The Phillips Curve and NAIRU}
The original Phillips curve suggested a stable trade-off between inflation
and unemployment. What historical episode challenged this relationship?
\begin{choices}
  \choice        The Great Depression, when both inflation and unemployment
                 were zero.
  \correctchoice Stagflation in the 1970s, when high inflation and high
                 unemployment occurred simultaneously.
  \choice        The introduction of the euro, which eliminated inflation
                 entirely.
  \choice        The 2008 financial crisis, which caused deflation and full
                 employment simultaneously.
\end{choices}
\begin{solution}
  The classic Phillips curve implied that policymakers could trade off lower
  unemployment for higher inflation. The 1970s \textbf{stagflation}
  episode --- driven partly by oil price shocks --- produced simultaneously
  high inflation \emph{and} high unemployment, contradicting the simple
  trade-off. This led to the concept of the \textbf{NAIRU} (Non-Accelerating
  Inflation Rate of Unemployment) and the recognition that supply shocks can
  shift the curve, making the inflation-unemployment relationship far less
  stable than originally believed.
\end{solution}
\end{question}

\begin{question}[type=mc]{The US Dollar's Global Role}
The script notes that the US dollar's importance in global finance goes
beyond its role as a reserve currency. What else supports its dominant
position?
\begin{choices}
  \choice        The US dollar is the only currency accepted in
                 international shipping contracts.
  \correctchoice The depth and liquidity of USD-denominated asset markets
                 (such as US Treasuries), which are widely held internationally
                 alongside dollars themselves.
  \choice        All countries are legally required to hold a minimum share
                 of USD in their reserves.
  \choice        The US dollar is the only currency accepted for oil trades
                 worldwide.
\end{choices}
\begin{solution}
  The dollar's dominance rests not only on central banks holding it as a
  \textbf{reserve currency}, but on the vast size and liquidity of
  \textbf{USD-denominated asset markets} --- above all US Treasury
  securities (USTs) --- which foreign investors, governments, and
  institutions hold extensively (alongside comparable assets like UK Gilts
  or Japanese government bonds). This depth of liquid, trusted assets is a
  key reason the dollar remains central to the global financial system, even
  as its dominance faces periodic challenges.
\end{solution}
\end{question}


%% ═════════════════════════════════════════════════════════════════════════════
\section{Chapter 10 — Macroeconomics in the Open Economy}
%% ═════════════════════════════════════════════════════════════════════════════

\begin{question}[type=mc]{Purchasing Power Parity}
The ``Big Mac Index'' is a popular illustration of Purchasing Power Parity
(PPP). What does PPP theory predict about exchange rates in the long run?
\begin{choices}
  \choice        Exchange rates should be fixed by central banks to prevent
                 any deviation from PPP.
  \correctchoice The exchange rate should adjust so that a given basket of
                 goods costs the same amount in different countries once
                 converted into a common currency.
  \choice        PPP applies equally well to tradable and non-tradable goods.
  \choice        PPP predicts that wages will equalise instantly across
                 countries.
\end{choices}
\begin{solution}
  PPP theory says that, in the long run and with free trade, the exchange
  rate should adjust so that the \textbf{same bundle of goods} costs the
  same amount everywhere once converted into a common currency --- money
  should have the same purchasing power at home and abroad. In practice, PPP
  holds much better for \textbf{tradable} goods (arbitraged across borders)
  than for \textbf{non-tradable} services like haircuts, which instead
  reflect local labour costs and productivity.
\end{solution}
\end{question}

\begin{question}[type=mc]{The Marshall-Lerner Condition}
The Marshall-Lerner condition, $M = |\eta_{IM}| + |\eta_{Ex}| > 1$, tells us
that:
\begin{choices}
  \choice        A currency depreciation always worsens the trade balance,
                 regardless of elasticities.
  \correctchoice A currency depreciation improves the trade balance only if
                 the combined absolute elasticities of import and export
                 demand exceed 1; if they are too low, depreciation may fail
                 to help or even hurt the trade balance.
  \choice        The condition only applies to fixed exchange rate regimes.
  \choice        A depreciation always leads to hyperinflation if $M > 1$.
\end{choices}
\begin{solution}
  If $M > 1$, import and export volumes respond strongly enough to a
  \textbf{depreciation} that the trade balance improves (the standard, and
  more common, case). If $M = 1$, depreciation has no effect on the trade
  balance; if $M < 1$, the exchange market is considered unstable and a
  depreciation can actually \textbf{worsen} the trade balance, because
  export and import volumes barely respond to the more favourable prices ---
  the only visible effect being higher import price inflation.
\end{solution}
\end{question}

\begin{question}[type=mc]{The Mundell-Fleming Model}
The Mundell-Fleming model extends the closed-economy IS-LM model by adding:
\begin{choices}
  \choice        A Phillips curve linking inflation to unemployment.
  \correctchoice A ZZ (Balance-of-Payments) curve, showing combinations of
                 interest rate and income that keep the balance of payments
                 in equilibrium, given the degree of international capital
                 mobility.
  \choice        A Laffer curve linking tax rates to government revenue.
  \choice        A production possibility frontier for the traded-goods
                 sector.
\end{choices}
\begin{solution}
  Mundell-Fleming adds a third curve, the \textbf{ZZ curve} (or BoP curve),
  to the IS-LM framework: all combinations of interest rate and output at
  which the balance of payments balances. If capital is perfectly mobile,
  the domestic interest rate is effectively pinned to the world interest
  rate. Whether fiscal or monetary policy is more effective then depends
  critically on whether the \textbf{exchange rate regime} is fixed or
  flexible --- under perfect capital mobility, monetary policy is powerful
  under flexible rates but powerless under fixed rates, and vice versa for
  fiscal policy.
\end{solution}
\end{question}

\end{document}
